\chapter{Resultados}

% --------------------------------------------------------
\section{Baseline ARIMA}
% --------------------------------------------------------

El modelo ARIMA(2,0,1), seleccionado automáticamente mediante criterio AIC, establece el piso de comparación.
La prueba ADF sobre el EMBI arroja estadístico $-4.336$ ($p = 0.0004$), confirmando estacionariedad en niveles.

\begin{table}[H]
\centering
\caption{Desempeño del baseline ARIMA}
\label{tab:arima}
\begin{tabular}{lcc}
\toprule
\textbf{Período} & \textbf{Esquema} & \textbf{RMSE (pb)} \\
\midrule
Validación & Rolling window (3 años) & 59.47 \\
Test        & Rolling window (3 años) & 107.48 \\
\bottomrule
\end{tabular}
\end{table}

% TODO: comentar el salto validación→test como evidencia de que ARIMA no captura shocks exógenos

% --------------------------------------------------------
\section{Selección de modelo ML: XGBoost vs Random Forest}
% --------------------------------------------------------

\begin{table}[H]
\centering
\caption{Comparación XGBoost vs Random Forest (CV, 10 folds)}
\label{tab:cv_comparison}
\begin{tabular}{lccc}
\toprule
\textbf{Modelo} & \textbf{RMSE CV medio (pb)} & \textbf{Std CV (pb)} & \textbf{p-valor Wilcoxon} \\
\midrule
Random Forest & 299.29 & $\pm$494.57 & \multirow{2}{*}{0.0273} \\
XGBoost       & 241.55 & $\pm$499.00 & \\
\bottomrule
\end{tabular}
\end{table}

% TODO: comentar el alto std como característica del EMBI ecuatoriano con picos extremos

Los hiperparámetros óptimos del XGBoost son: \texttt{n\_estimators=607}, \texttt{learning\_rate=0.093},
\texttt{max\_depth=4}, \texttt{subsample=0.848}, \texttt{colsample\_bytree=0.980},
\texttt{reg\_alpha=0.662}, \texttt{reg\_lambda=1.459}, \texttt{gamma=0.735}.

% --------------------------------------------------------
\section{Desempeño en test: XGBoost completo}
% --------------------------------------------------------

\begin{table}[H]
\centering
\caption{Métricas del XGBoost completo en test}
\label{tab:xgb_test}
\begin{tabular}{lc}
\toprule
\textbf{Métrica} & \textbf{Valor} \\
\midrule
RMSE  & 64.16 pb \\
MAE   & 34.90 pb \\
R$^2$ & 0.9555   \\
Mejora vs ARIMA & $-$40.3\% \\
\bottomrule
\end{tabular}
\end{table}

% --------------------------------------------------------
\section{Experimento de ablación}
% --------------------------------------------------------

% TODO: introducir los 4 grupos y su justificación

\begin{table}[H]
\centering
\caption{Resultados del experimento de ablación por grupo de variables}
\label{tab:ablation}
\begin{tabular}{lcccc}
\toprule
\textbf{Grupo} & \textbf{Features NLP} & \textbf{RMSE (pb)} & \textbf{MAE (pb)} & \textbf{R$^2$} \\
\midrule
ARIMA baseline      & ---  & 107.48 & ---   & ---    \\
A — Solo AR         & 0    & 51.31  & 27.15 & 0.9715 \\
B — AR + Macro      & 0    & 60.07  & 41.12 & 0.9610 \\
C — AR + NLP        & 4    & \textbf{49.01}  & \textbf{24.93} & \textbf{0.9740} \\
D — AR + Macro + NLP& 4    & 57.48  & 37.70 & 0.9643 \\
\bottomrule
\end{tabular}
\end{table}

% TODO: comentar hallazgos principales:
%   1. Grupo B peor que Grupo A — variables macro locales generan ruido
%   2. Grupo C mejor — NLP añade señal complementaria a los lags AR
%   3. Grupo D no es el mejor — combinación macro+NLP diluye el beneficio del NLP

% Wilcoxon tests
% TODO: reportar p-valores (A vs C: 0.1250, B vs D: 0.3125) y discutir bajo poder con 5 folds

% Figuras SHAP
\begin{figure}[H]
\centering
\includegraphics[width=0.85\textwidth]{shap_beeswarm.png}
\caption{SHAP beeswarm — importancia y dirección de las features en el modelo completo}
\label{fig:shap_beeswarm}
\end{figure}

\begin{figure}[H]
\centering
\includegraphics[width=0.85\textwidth]{shap_categorias.png}
\caption{Importancia SHAP agregada por categoría de variable}
\label{fig:shap_categorias}
\end{figure}

% --------------------------------------------------------
\section{Causalidad de Granger}
% --------------------------------------------------------

% TODO: tabla de resultados Granger (dirección, p-valor, lags significativos)
% Figura
\begin{figure}[H]
\centering
\includegraphics[width=0.75\textwidth]{granger_heatmap.png}
\caption{Heatmap de causalidad de Granger (lags 1--10 días hábiles)}
\label{fig:granger}
\end{figure}

% TODO: comentar: sin causalidad NLP→EMBI; causalidad inversa EMBI→tono (lags 1-4)

% --------------------------------------------------------
\section{Análisis de regímenes estructurales}
% --------------------------------------------------------

\subsection{Tests de quiebre estructural}

\begin{table}[H]
\centering
\caption{Resultados de los tests de quiebre estructural}
\label{tab:chow}
\begin{tabular}{lccccc}
\toprule
\textbf{Test} & \textbf{Fecha} & \textbf{F-stat} & \textbf{p-valor} & \textbf{n$_1$} & \textbf{n$_2$} \\
\midrule
Chow  & 2020-03-01 & 5.52  & 0.0002 & 1.587 & 1.457 \\
Chow  & 2022-01-01 & 11.61 & 0.0000 & 2.048 & 996   \\
Quandt-Andrews & 2020-03-25 & 53.58 & 0.0000 & --- & --- \\
CUSUM & --- & 0.848 & 0.469 & \multicolumn{2}{c}{No significativo} \\
\bottomrule
\end{tabular}
\end{table}

\begin{figure}[H]
\centering
\includegraphics[width=0.85\textwidth]{chow_cusum_analysis.png}
\caption{CUSUM sobre residuos OLS AR(3) con bandas de confianza al 95\%}
\label{fig:cusum}
\end{figure}

\subsection{Importancia SHAP por régimen}

\begin{figure}[H]
\centering
\includegraphics[width=\textwidth]{shap_regimenes_heatmap.png}
\caption{Heatmap de importancia SHAP (top-15 features) por régimen temporal}
\label{fig:shap_regimenes}
\end{figure}

\begin{figure}[H]
\centering
\includegraphics[width=0.75\textwidth]{rmse_por_regimen.png}
\caption{RMSE por subperíodo: pre-COVID, COVID, post-COVID}
\label{fig:rmse_regimen}
\end{figure}

% TODO: comentar los tres hallazgos del régimen:
%   1. Pre-COVID: exogeneidad total (variables externas + AR)
%   2. COVID: desempleo local irrumpe (SHAP=125.98)
%   3. Post-COVID: dominancia Fed + DXY en economía dolarizada

% --------------------------------------------------------
\section{Autoencoder NLP}
% --------------------------------------------------------

\begin{figure}[H]
\centering
\includegraphics[width=0.7\textwidth]{autoencoder_loss_curve.png}
\caption{Curva de pérdida del autoencoder por época (train vs validación)}
\label{fig:ae_loss}
\end{figure}

\begin{table}[H]
\centering
\caption{Comparativa de representaciones NLP en el modelo XGBoost}
\label{tab:autoencoder}
\begin{tabular}{lcccc}
\toprule
\textbf{Representación NLP} & \textbf{Features} & \textbf{RMSE (pb)} & \textbf{MAE (pb)} & \textbf{R$^2$} \\
\midrule
Original (4 rolling 30d) & 4 & 59.29 & 39.37 & 0.9620 \\
PCA-3 (reducción lineal) & 3 & 60.11 & 41.14 & 0.9609 \\
Autoencoder-3 (no lineal)& 3 & 63.41 & 44.79 & 0.9565 \\
\bottomrule
\end{tabular}
\end{table}

\begin{figure}[H]
\centering
\includegraphics[width=0.75\textwidth]{nlp_latentes_correlacion.png}
\caption{Correlación entre dimensiones latentes del autoencoder y variables NLP originales}
\label{fig:ae_latentes}
\end{figure}

% TODO: comentar que la estructura NLP es esencialmente lineal (PCA 77.7% varianza en 3 comp)
% y que comprimir de 4→3 ya pierde señal predictiva marginal
